(1) 表面电荷随球转动的切向速度 \(v_\phi=\Omega a\sin\theta\)，故 \(\boldsymbol K=\sigma\Omega a\sin\theta\,\boldsymbol e_\phi\)。用表面电流磁偶极定义
\(\boldsymbol m=\tfrac12\int_{S_a}\boldsymbol r\times\boldsymbol K\,dS\)，其 \(z\) 分量为
\[
m_z=\tfrac12\sigma\Omega a^2\int_{S_a}\sin^2\theta\,dS
=\tfrac12\sigma\Omega a^4(8\pi/3)=Q\Omega a^2/3.
\]
方向沿正转轴，任何把旋转球壳当单圆环而得 \(Q\Omega a^2/2\) 的做法都忽略了纬度权重。

(2) 旋转表面电流等价于均匀磁化强度 \(\boldsymbol M=\sigma\Omega a\,\boldsymbol e_z\) 的球面束缚电流 \(\boldsymbol K=\boldsymbol M\times\boldsymbol n\)。均匀磁化球的内部退磁场 \(\boldsymbol H_{\rm in}=-\boldsymbol M/3\)，故
\[
\boldsymbol B_{\rm in}=\frac23\mu_0\boldsymbol M=\frac{\mu_0Q\Omega}{6\pi a}\boldsymbol e_z.
\]
球外场为 \(\boldsymbol B_{\rm out}=\frac{\mu_0}{4\pi r^3}[3(\boldsymbol m\cdot\hat r)\hat r-\boldsymbol m]\)。在表面 \(B_r\) 连续，切向跳跃 \(\boldsymbol n\times(\boldsymbol H_{\rm out}-\boldsymbol H_{\rm in})=\boldsymbol K\)，可代入球坐标的 \(\sin\theta\) 分量严格核验。

(3) 用 \(B_{\rm in}=\mu_0m/(2\pi a^3)\)，内部能量
\[
U_{\rm in}=\frac{B_{\rm in}^2}{2\mu_0}\frac{4\pi a^3}{3}
=\frac{\mu_0m^2}{6\pi a^3}.
\]
外部偶极 \(B^2=\frac{\mu_0^2m^2}{16\pi^2r^6}(1+3\cos^2\theta)\)，角积分给 \(8\pi\)，径向积分 \(\int_a^\infty r^{-4}dr=1/(3a^3)\)，所以 \(U_{\rm out}=\mu_0m^2/(12\pi a^3)\)。总和
\(\boxed{U_B=\mu_0m^2/(4\pi a^3)}\)；独立取 \(A_\phi(a)=\mu_0m\sin\theta/(4\pi a^2)\) 求半倍表面 \(KA\) 积分也给同一结果。

(4) 固定 \(Q,\Omega\) 时 \(m\propto a^2\)，从而 \(U_B\propto a\)；此极限在形式上趋零，但不能同时维持球壳的实际材料、机械稳定与点电荷自能有限。若固定磁矩 \(m\) 反而要 \(\Omega\propto a^{-2}\)，违反 \(\Omega a\ll c\)，并使场能按 \(a^{-3}\) 发散。静态 \(\boldsymbol J\) 与磁场没有辐射、涡流或黏性阻力机制，故单靠上述无阻模型无法推断转动阻矩；启停过程还涉及时变感生场。
